\section*{Metodología}

El análisis presentado en este artículo no parte de un único enfoque, porque el propio tema —la arquitectura por capas— tampoco puede estudiarse desde una sola perspectiva. La metodología utilizada combina una revisión documental amplia con un ejercicio interpretativo que intenta poner en relación trabajos que, a primera vista, parecen hablar de asuntos distintos. Más que buscar una síntesis perfecta, el propósito fue entender qué tensiones y qué coincidencias aparecen cuando se observa el patrón desde contextos variados.

El primer paso consistió en revisar investigaciones que abarcan diferentes momentos y motivaciones: estudios que examinan la erosión arquitectónica, propuestas para visualizar dependencias, análisis de rendimiento basados en modelos probabilísticos, evaluaciones empíricas en microservicios y discusiones sobre cómo adaptar las capas a tecnologías emergentes. No se trató de clasificarlas de forma rígida, sino de identificar qué preguntas están intentando responder y qué tipo de problemas reconocen en la práctica.

Posteriormente, se organizaron los hallazgos según el tipo de aporte que ofrecen. Algunos trabajos describen fenómenos recurrentes en sistemas grandes; otros presentan herramientas para interpretar estructuras que ya no son evidentes; y también existen propuestas que buscan flexibilizar o extender el modelo clásico. Esta agrupación permitió observar patrones comunes, incluso entre investigaciones que no se citan entre sí.

Además de la revisión teórica, se incorporó un análisis del caso de un sistema real: una plataforma de carnetización digital que reúne registro de personas, autenticación, gestión de eventos y control de asistencia. El objetivo no fue convertir este sistema en un estudio de caso formal, sino usarlo como una referencia práctica para evaluar hasta qué punto los principios encontrados en la literatura se alinean o chocan con lo que ocurre en un proyecto en funcionamiento. Este caso sirvió como una forma de comprobar qué partes del modelo en capas resultan útiles y cuáles requieren adaptaciones para sostener la operación diaria.

En conjunto, la metodología siguió un enfoque cualitativo y comparativo, buscó construir una mirada más amplia sobre el enfoque arquitectónico en capas. El resultado final surge de la interacción entre los estudios consultados y la experiencia práctica observada en un sistema contemporáneo que enfrenta desafíos típicos del desarrollo moderno.
